
\documentclass{article}
\usepackage{colortbl}
\usepackage{makecell}
\usepackage{multirow}
\usepackage{supertabular}

\begin{document}

\newcounter{utterance}

\centering \large Interaction Transcript for game `clean\_up', experiment `1\_medium\_7obj\_es', episode 3 with Apertus{-}v1.5{-}8B.
\vspace{24pt}

{ \footnotesize  \setcounter{utterance}{1}
\setlength{\tabcolsep}{0pt}
\begin{supertabular}{c@{$\;$}|p{.15\linewidth}@{}p{.15\linewidth}p{.15\linewidth}p{.15\linewidth}p{.15\linewidth}p{.15\linewidth}}
    \# & $\;$A & \multicolumn{4}{c}{Game Master} & $\;\:$B\\
    \hline

    \theutterance \stepcounter{utterance}  
    & & \multicolumn{4}{p{0.6\linewidth}}{
        \cellcolor[rgb]{0.9,0.9,0.9}{
            \makecell[{{p{\linewidth}}}]{
                \texttt{\tiny{[P1$\langle$GM]}}
                \texttt{Soy tu maestro de juego, y estás jugando un juego colaborativo con la siguiente cuadrícula como tablero de juego:} \\
\\ 
\texttt{```} \\
\texttt{ 1234567} \\
\texttt{╔═╤═════╗} \\
\texttt{║◌│◌A◌◌◌║ 1} \\
\texttt{╟─┤◌T◌◌◌║ 2} \\
\texttt{║O└─┐◌◌◌║ 3} \\
\texttt{║◌◌◌├─┐◌║ 4} \\
\texttt{║◌◌S│◌│◌║ 5} \\
\texttt{╟─┬─┴─┼─╢ 6} \\
\texttt{║P│◌DM│◌║ 7} \\
\texttt{╚═╧═══╧═╝} \\
\\ 
\texttt{```} \\
\\ 
\texttt{* El borde superior muestra coordenadas x que aumentan hacia la derecha, y el borde derecho coordenadas y que aumentan hacia abajo.} \\
\texttt{* Los siguientes objetos están colocados aleatoriamente en tu cuadrícula: 'P', 'O', 'T', 'S', 'D', 'A', 'M'.} \\
\\ 
\\ 
\texttt{El otro jugador ve una variación del tablero de juego, donde los objetos están colocados en diferentes ubicaciones aleatorias. No puedes ver el tablero del otro jugador, y él no puede ver el tuyo.} \\
\\ 
\\ 
\texttt{**Objetivo:**} \\
\\ 
\texttt{Ambos jugadores necesitan mover los objetos en sus respectivos fondos de manera que los objetos idénticos terminen en las mismas coordenadas. Tienes que comunicarte con el otro jugador para acordar una configuración objetivo común.} \\
\\ 
\\ 
\texttt{**Reglas:**} \\
\\ 
\texttt{* En cada turno, puedes enviar exactamente uno de los dos comandos siguientes:} \\
\texttt{ 1. `DECIR: <MENSAJE>`: para enviar un mensaje (todo hasta el siguiente salto de línea) al otro jugador. Se lo reenviaré a tu compañero.} \\
\texttt{2. `MOVERR: <OBJETO>, (<X>, <Y>)`: para mover un objeto a una nueva posición, donde `<X>` es la columna y `<Y>` es la fila. Te informaré si tu movimiento fue exitoso o no.} \\
\texttt{* Si no sigues el formato, o envías varios comandos a la vez, tendré que penalizarte.} \\
\texttt{* Si ambos jugadores acumulan más de 16 penalizaciones, ambos pierden el juego.} \\
\texttt{* ¡Es vital que te comuniques con el otro jugador respecto a tu estado objetivo! ¡La *única* forma de transmitir tu estrategia al otro jugador es usando el comando `DECIR: <MENSAJE>`!} \\
\\ 
\\ 
\texttt{**Mover objetos**} \\
\\ 
\texttt{* Solo puedes mover objetos a celdas dentro de los límites de la cuadrícula. La celda de destino debe estar vacía, es decir, solo debe contener el símbolo '◌'.} \\
\texttt{* Si intentas mover un objeto a un lugar que no está vacío, o intentas moverlo fuera de la cuadrícula, tendré que penalizarte. Tendrás otro intento.} \\
\texttt{* ¡Antes de hacer un movimiento, verifica dos veces que el lugar de destino esté vacío, y no contenga ninguna letra, marco o línea!} \\
\\ 
\\ 
\texttt{**Fin del juego**} \\
\\ 
\texttt{Si crees que alcanzaste el objetivo de alinear todos los objetos, puedes pedirle al otro jugador que termine el juego enviando `DECIR: ¿terminado?`. Si el otro jugador te pide que termines el juego, y respondes con `DECIR: ¡terminado!`, el juego terminará.} \\
\\ 
\texttt{Ambos jugadores ganan si el juego termina dentro de 28 rondas, donde una ronda se define como dos jugadores enviando cada uno un comando válido.} \\
\\ 
\texttt{**Puntuación:**} \\
\\ 
\texttt{Cuanto más cerca estén los objetos idénticos en ambos tableros de juego, más puntos obtendrás. Las penalizaciones reducen tus puntos. ¿Puedes batir el récord?} \\
\\ 
\\ 
\texttt{¡Por favor, envía un mensaje al otro jugador para iniciar el juego!} \\
            }
        }
    }
    & \\ \\

    \theutterance \stepcounter{utterance}  
    & \multicolumn{4}{p{0.6\linewidth}}{
        \cellcolor[rgb]{0.9,0.9,0.9}{
            \makecell[{{p{\linewidth}}}]{
                \texttt{\tiny{[P1$\rangle$GM]}}
                \texttt{DECIR: Hola, ¿qué objetos tienes en tu tablero? Necesito saber qué debo buscar para coordinar nuestra estrategia.} \\
            }
        }
    }
    & & \\ \\

    \theutterance \stepcounter{utterance}  
    & & & \multicolumn{2}{p{0.3\linewidth}}{
        \cellcolor[rgb]{0.9,0.9,0.9}{
            \makecell[{{p{\linewidth}}}]{
                \texttt{\tiny{[GM$|$GM]}}
                \texttt{DECIR: Hola, ¿qué objetos tienes en tu tablero? Necesito saber qué debo buscar para coordinar nuestra estrategia.} \\
            }
        }
    }
    & & \\ \\

    \theutterance \stepcounter{utterance}  
    & & \multicolumn{4}{p{0.6\linewidth}}{
        \cellcolor[rgb]{0.9,0.9,0.9}{
            \makecell[{{p{\linewidth}}}]{
                \texttt{\tiny{[GM$\rangle$P2]}}
                \texttt{Soy tu maestro de juego, y estás jugando un juego colaborativo con la siguiente cuadrícula como tablero de juego:} \\
\\ 
\texttt{```} \\
\texttt{ 1234567} \\
\texttt{╔═╤═════╗} \\
\texttt{║P│S◌D◌◌║ 1} \\
\texttt{╟─┤◌◌◌◌◌║ 2} \\
\texttt{║◌└─┐◌A◌║ 3} \\
\texttt{║◌◌M├─┐◌║ 4} \\
\texttt{║◌T◌│◌│O║ 5} \\
\texttt{╟─┬─┴─┼─╢ 6} \\
\texttt{║◌│◌◌◌│◌║ 7} \\
\texttt{╚═╧═══╧═╝} \\
\\ 
\texttt{```} \\
\\ 
\texttt{* El borde superior muestra coordenadas x que aumentan hacia la derecha, y el borde derecho coordenadas y que aumentan hacia abajo.} \\
\texttt{* Los siguientes objetos están colocados aleatoriamente en tu cuadrícula: 'P', 'O', 'T', 'S', 'D', 'A', 'M'.} \\
\\ 
\\ 
\texttt{El otro jugador ve una variación del tablero de juego, donde los objetos están colocados en diferentes ubicaciones aleatorias. No puedes ver el tablero del otro jugador, y él no puede ver el tuyo.} \\
\\ 
\\ 
\texttt{**Objetivo:**} \\
\\ 
\texttt{Ambos jugadores necesitan mover los objetos en sus respectivos fondos de manera que los objetos idénticos terminen en las mismas coordenadas. Tienes que comunicarte con el otro jugador para acordar una configuración objetivo común.} \\
\\ 
\\ 
\texttt{**Reglas:**} \\
\\ 
\texttt{* En cada turno, puedes enviar exactamente uno de los dos comandos siguientes:} \\
\texttt{ 1. `DECIR: <MENSAJE>`: para enviar un mensaje (todo hasta el siguiente salto de línea) al otro jugador. Se lo reenviaré a tu compañero.} \\
\texttt{2. `MOVERR: <OBJETO>, (<X>, <Y>)`: para mover un objeto a una nueva posición, donde `<X>` es la columna y `<Y>` es la fila. Te informaré si tu movimiento fue exitoso o no.} \\
\texttt{* Si no sigues el formato, o envías varios comandos a la vez, tendré que penalizarte.} \\
\texttt{* Si ambos jugadores acumulan más de 16 penalizaciones, ambos pierden el juego.} \\
\texttt{* ¡Es vital que te comuniques con el otro jugador respecto a tu estado objetivo! ¡La *única* forma de transmitir tu estrategia al otro jugador es usando el comando `DECIR: <MENSAJE>`!} \\
\\ 
\\ 
\texttt{**Mover objetos**} \\
\\ 
\texttt{* Solo puedes mover objetos a celdas dentro de los límites de la cuadrícula. La celda de destino debe estar vacía, es decir, solo debe contener el símbolo '◌'.} \\
\texttt{* Si intentas mover un objeto a un lugar que no está vacío, o intentas moverlo fuera de la cuadrícula, tendré que penalizarte. Tendrás otro intento.} \\
\texttt{* ¡Antes de hacer un movimiento, verifica dos veces que el lugar de destino esté vacío, y no contenga ninguna letra, marco o línea!} \\
\\ 
\\ 
\texttt{**Fin del juego**} \\
\\ 
\texttt{Si crees que alcanzaste el objetivo de alinear todos los objetos, puedes pedirle al otro jugador que termine el juego enviando `DECIR: ¿terminado?`. Si el otro jugador te pide que termines el juego, y respondes con `DECIR: ¡terminado!`, el juego terminará.} \\
\\ 
\texttt{Ambos jugadores ganan si el juego termina dentro de 28 rondas, donde una ronda se define como dos jugadores enviando cada uno un comando válido.} \\
\\ 
\texttt{**Puntuación:**} \\
\\ 
\texttt{Cuanto más cerca estén los objetos idénticos en ambos tableros de juego, más puntos obtendrás. Las penalizaciones reducen tus puntos. ¿Puedes batir el récord?} \\
\\ 
\texttt{El otro jugador inició el juego enviando este mensaje:} \\
\\ 
\texttt{"Hola, ¿qué objetos tienes en tu tablero? Necesito saber qué debo buscar para coordinar nuestra estrategia."} \\
\\ 
\texttt{¿Cuál es tu primer comando?} \\
            }
        }
    }
    & \\ \\

    \theutterance \stepcounter{utterance}  
    & & & \multicolumn{4}{p{0.6\linewidth}}{
        \cellcolor[rgb]{0.9,0.9,0.9}{
            \makecell[{{p{\linewidth}}}]{
                \texttt{\tiny{[GM$\langle$P2]}}
                \texttt{DECIR: Tengo los siguientes objetos en mi tablero: 'P', 'O', 'T', 'S', 'D', 'A', 'M'. Por favor, confirma si esta información es correcta antes de proceder.} \\
            }
        }
    }
     \\ \\

    \theutterance \stepcounter{utterance}  
    & & & \multicolumn{2}{p{0.3\linewidth}}{
        \cellcolor[rgb]{0.9,0.9,0.9}{
            \makecell[{{p{\linewidth}}}]{
                \texttt{\tiny{[GM$|$GM]}}
                \texttt{DECIR: Tengo los siguientes objetos en mi tablero: 'P', 'O', 'T', 'S', 'D', 'A', 'M'. Por favor, confirma si esta información es correcta antes de proceder.} \\
            }
        }
    }
    & & \\ \\

    \theutterance \stepcounter{utterance}  
    & & \multicolumn{4}{p{0.6\linewidth}}{
        \cellcolor[rgb]{0.9,0.9,0.9}{
            \makecell[{{p{\linewidth}}}]{
                \texttt{\tiny{[P1$\langle$GM]}}
                \texttt{Tu mensaje ha sido transmitido al otro jugador.} \\
\\ 
\\ 
\texttt{Actualmente estás jugando la ronda 1 de un máximo de 28.} \\
\texttt{Han acumulado colectivamente 0 de 16 penalizaciones.} \\
\texttt{El otro jugador envió este mensaje:} \\
\\ 
\\ 
\texttt{"Tengo los siguientes objetos en mi tablero: 'P', 'O', 'T', 'S', 'D', 'A', 'M'. Por favor, confirma si esta información es correcta antes de proceder."} \\
\\ 
\\ 
\texttt{¿Cuál es tu próximo comando?} \\
            }
        }
    }
    & \\ \\

    \theutterance \stepcounter{utterance}  
    & \multicolumn{4}{p{0.6\linewidth}}{
        \cellcolor[rgb]{0.9,0.9,0.9}{
            \makecell[{{p{\linewidth}}}]{
                \texttt{\tiny{[P1$\rangle$GM]}}
                \texttt{DECIR: Confirmo, mi tablero también tiene los objetos: 'P', 'O', 'T', 'S', 'D', 'A', 'M'. Estoy listo para coordinar movimientos y alinear todos los objetos en las mismas posiciones.} \\
            }
        }
    }
    & & \\ \\

    \theutterance \stepcounter{utterance}  
    & & & \multicolumn{2}{p{0.3\linewidth}}{
        \cellcolor[rgb]{0.9,0.9,0.9}{
            \makecell[{{p{\linewidth}}}]{
                \texttt{\tiny{[GM$|$GM]}}
                \texttt{DECIR: Confirmo, mi tablero también tiene los objetos: 'P', 'O', 'T', 'S', 'D', 'A', 'M'. Estoy listo para coordinar movimientos y alinear todos los objetos en las mismas posiciones.} \\
            }
        }
    }
    & & \\ \\

    \theutterance \stepcounter{utterance}  
    & & \multicolumn{4}{p{0.6\linewidth}}{
        \cellcolor[rgb]{0.9,0.9,0.9}{
            \makecell[{{p{\linewidth}}}]{
                \texttt{\tiny{[GM$\rangle$P2]}}
                \texttt{Tu mensaje ha sido transmitido al otro jugador.} \\
\\ 
\\ 
\texttt{Actualmente estás jugando la ronda 2 de un máximo de 28.} \\
\texttt{Han acumulado colectivamente 0 de 16 penalizaciones.} \\
\texttt{El otro jugador envió este mensaje:} \\
\\ 
\\ 
\texttt{"Confirmo, mi tablero también tiene los objetos: 'P', 'O', 'T', 'S', 'D', 'A', 'M'. Estoy listo para coordinar movimientos y alinear todos los objetos en las mismas posiciones."} \\
\\ 
\\ 
\texttt{¿Cuál es tu próximo comando?} \\
            }
        }
    }
    & \\ \\

    \theutterance \stepcounter{utterance}  
    & & & \multicolumn{4}{p{0.6\linewidth}}{
        \cellcolor[rgb]{0.9,0.9,0.9}{
            \makecell[{{p{\linewidth}}}]{
                \texttt{\tiny{[GM$\langle$P2]}}
                \texttt{DECIR: De acuerdo, confirmo que ambos tableros contienen los mismos objetos: 'P', 'O', 'T', 'S', 'D', 'A', 'M'. Para proceder, necesito que me proporciones tu estrategia para alinear los objetos en posiciones comunes. ¿Podrías describir brevemente cómo planeas mover los objetos para lograr este objetivo? Estoy listo para coordinar movimientos y compartir mi propia estrategia también.} \\
            }
        }
    }
     \\ \\

    \theutterance \stepcounter{utterance}  
    & & & \multicolumn{2}{p{0.3\linewidth}}{
        \cellcolor[rgb]{0.9,0.9,0.9}{
            \makecell[{{p{\linewidth}}}]{
                \texttt{\tiny{[GM$|$GM]}}
                \texttt{DECIR: De acuerdo, confirmo que ambos tableros contienen los mismos objetos: 'P', 'O', 'T', 'S', 'D', 'A', 'M'. Para proceder, necesito que me proporciones tu estrategia para alinear los objetos en posiciones comunes. ¿Podrías describir brevemente cómo planeas mover los objetos para lograr este objetivo? Estoy listo para coordinar movimientos y compartir mi propia estrategia también.} \\
            }
        }
    }
    & & \\ \\

    \theutterance \stepcounter{utterance}  
    & & \multicolumn{4}{p{0.6\linewidth}}{
        \cellcolor[rgb]{0.9,0.9,0.9}{
            \makecell[{{p{\linewidth}}}]{
                \texttt{\tiny{[P1$\langle$GM]}}
                \texttt{Tu mensaje ha sido transmitido al otro jugador.} \\
\\ 
\\ 
\texttt{Actualmente estás jugando la ronda 2 de un máximo de 28.} \\
\texttt{Han acumulado colectivamente 0 de 16 penalizaciones.} \\
\texttt{El otro jugador envió este mensaje:} \\
\\ 
\\ 
\texttt{"De acuerdo, confirmo que ambos tableros contienen los mismos objetos: 'P', 'O', 'T', 'S', 'D', 'A', 'M'. Para proceder, necesito que me proporciones tu estrategia para alinear los objetos en posiciones comunes. ¿Podrías describir brevemente cómo planeas mover los objetos para lograr este objetivo? Estoy listo para coordinar movimientos y compartir mi propia estrategia también."} \\
\\ 
\\ 
\texttt{¿Cuál es tu próximo comando?} \\
            }
        }
    }
    & \\ \\

    \theutterance \stepcounter{utterance}  
    & \multicolumn{4}{p{0.6\linewidth}}{
        \cellcolor[rgb]{0.9,0.9,0.9}{
            \makecell[{{p{\linewidth}}}]{
                \texttt{\tiny{[P1$\rangle$GM]}}
                \texttt{DECIR: En mi estrategia, primero identificaré las posiciones de cada objeto en ambos tableros. Luego, intentaré mover objetos uno por uno, verificando siempre que la celda de destino esté vacía antes de cada movimiento. Comenzaré con los objetos que están más cerca de su posición objetivo ideal. Para la coordinación, voy a describir los movimientos iniciales que propongo y estaré atento a tu respuesta para ajustar la estrategia. ¿Podrías describir brevemente tu estrategia y las posiciones que tienes en mente para cada objeto?} \\
            }
        }
    }
    & & \\ \\

    \theutterance \stepcounter{utterance}  
    & & & \multicolumn{2}{p{0.3\linewidth}}{
        \cellcolor[rgb]{0.9,0.9,0.9}{
            \makecell[{{p{\linewidth}}}]{
                \texttt{\tiny{[GM$|$GM]}}
                \texttt{DECIR: En mi estrategia, primero identificaré las posiciones de cada objeto en ambos tableros. Luego, intentaré mover objetos uno por uno, verificando siempre que la celda de destino esté vacía antes de cada movimiento. Comenzaré con los objetos que están más cerca de su posición objetivo ideal. Para la coordinación, voy a describir los movimientos iniciales que propongo y estaré atento a tu respuesta para ajustar la estrategia. ¿Podrías describir brevemente tu estrategia y las posiciones que tienes en mente para cada objeto?} \\
            }
        }
    }
    & & \\ \\

    \theutterance \stepcounter{utterance}  
    & & \multicolumn{4}{p{0.6\linewidth}}{
        \cellcolor[rgb]{0.9,0.9,0.9}{
            \makecell[{{p{\linewidth}}}]{
                \texttt{\tiny{[GM$\rangle$P2]}}
                \texttt{Tu mensaje ha sido transmitido al otro jugador.} \\
\\ 
\\ 
\texttt{Actualmente estás jugando la ronda 3 de un máximo de 28.} \\
\texttt{Han acumulado colectivamente 0 de 16 penalizaciones.} \\
\texttt{El otro jugador envió este mensaje:} \\
\\ 
\\ 
\texttt{"En mi estrategia, primero identificaré las posiciones de cada objeto en ambos tableros. Luego, intentaré mover objetos uno por uno, verificando siempre que la celda de destino esté vacía antes de cada movimiento. Comenzaré con los objetos que están más cerca de su posición objetivo ideal. Para la coordinación, voy a describir los movimientos iniciales que propongo y estaré atento a tu respuesta para ajustar la estrategia. ¿Podrías describir brevemente tu estrategia y las posiciones que tienes en mente para cada objeto?"} \\
\\ 
\\ 
\texttt{¿Cuál es tu próximo comando?} \\
            }
        }
    }
    & \\ \\

    \theutterance \stepcounter{utterance}  
    & & & \multicolumn{4}{p{0.6\linewidth}}{
        \cellcolor[rgb]{0.9,0.9,0.9}{
            \makecell[{{p{\linewidth}}}]{
                \texttt{\tiny{[GM$\langle$P2]}}
                \texttt{DECIR: Entiendo tu enfoque. Mi estrategia será similar: identificar las posiciones actuales de cada objeto, verificar celdas vacías antes de mover, y priorizar objetos cercanos a su objetivo. Respecto a las posiciones actual de mis objetos en mi tablero, son: 'P' en (1,1), 'O' en (5,7), 'T' en (4,1), 'S' en (1,2), 'D' en (1,3), 'A' en (3,4), 'M' en (4,1). Estoy dispuesto a compartir detalles específicos de mi estrategia para coordinar movilizaciones, por ejemplo, nombres de movimientos propuestos o criterios de priorización, para que podamos adaptar acciones conjuntas y avanzar eficientemente. ¿Cuál es tu propuesta inicial para mover uno de los objetos, y cómo crees que deberíamos proceder para optimizar nuestro avance?} \\
            }
        }
    }
     \\ \\

    \theutterance \stepcounter{utterance}  
    & & & \multicolumn{2}{p{0.3\linewidth}}{
        \cellcolor[rgb]{0.9,0.9,0.9}{
            \makecell[{{p{\linewidth}}}]{
                \texttt{\tiny{[GM$|$GM]}}
                \texttt{DECIR: Entiendo tu enfoque. Mi estrategia será similar: identificar las posiciones actuales de cada objeto, verificar celdas vacías antes de mover, y priorizar objetos cercanos a su objetivo. Respecto a las posiciones actual de mis objetos en mi tablero, son: 'P' en (1,1), 'O' en (5,7), 'T' en (4,1), 'S' en (1,2), 'D' en (1,3), 'A' en (3,4), 'M' en (4,1). Estoy dispuesto a compartir detalles específicos de mi estrategia para coordinar movilizaciones, por ejemplo, nombres de movimientos propuestos o criterios de priorización, para que podamos adaptar acciones conjuntas y avanzar eficientemente. ¿Cuál es tu propuesta inicial para mover uno de los objetos, y cómo crees que deberíamos proceder para optimizar nuestro avance?} \\
            }
        }
    }
    & & \\ \\

    \theutterance \stepcounter{utterance}  
    & & \multicolumn{4}{p{0.6\linewidth}}{
        \cellcolor[rgb]{0.9,0.9,0.9}{
            \makecell[{{p{\linewidth}}}]{
                \texttt{\tiny{[P1$\langle$GM]}}
                \texttt{Tu mensaje ha sido transmitido al otro jugador.} \\
\\ 
\\ 
\texttt{Actualmente estás jugando la ronda 3 de un máximo de 28.} \\
\texttt{Han acumulado colectivamente 0 de 16 penalizaciones.} \\
\texttt{El otro jugador envió este mensaje:} \\
\\ 
\\ 
\texttt{"Entiendo tu enfoque. Mi estrategia será similar: identificar las posiciones actuales de cada objeto, verificar celdas vacías antes de mover, y priorizar objetos cercanos a su objetivo. Respecto a las posiciones actual de mis objetos en mi tablero, son: 'P' en (1,1), 'O' en (5,7), 'T' en (4,1), 'S' en (1,2), 'D' en (1,3), 'A' en (3,4), 'M' en (4,1). Estoy dispuesto a compartir detalles específicos de mi estrategia para coordinar movilizaciones, por ejemplo, nombres de movimientos propuestos o criterios de priorización, para que podamos adaptar acciones conjuntas y avanzar eficientemente. ¿Cuál es tu propuesta inicial para mover uno de los objetos, y cómo crees que deberíamos proceder para optimizar nuestro avance?"} \\
\\ 
\\ 
\texttt{¿Cuál es tu próximo comando?} \\
            }
        }
    }
    & \\ \\

    \theutterance \stepcounter{utterance}  
    & \multicolumn{4}{p{0.6\linewidth}}{
        \cellcolor[rgb]{0.9,0.9,0.9}{
            \makecell[{{p{\linewidth}}}]{
                \texttt{\tiny{[P1$\rangle$GM]}}
                \texttt{DECIR: Agradezco que compartas tu estrategia y posiciones. Para coordinar eficazmente, propongo comenzar con el objeto 'P', que se encuentra en (1,1). Observo que en mi tablero también tengo 'P' en (1,1). Si ambos estamos de acuerdo, podríamos intentar mover 'P' a una posición vacía que sea la más cercana para ambos tableros. ¿Podrías confirmar si, en tu tablero, hay una celda vacía cercana a (1,1) que ambos podríamos considerar como un punto de partida común para el movimiento de 'P'? Esta verificación es crucial para evitar penalizaciones en el primer movimiento.} \\
            }
        }
    }
    & & \\ \\

    \theutterance \stepcounter{utterance}  
    & & & \multicolumn{2}{p{0.3\linewidth}}{
        \cellcolor[rgb]{0.9,0.9,0.9}{
            \makecell[{{p{\linewidth}}}]{
                \texttt{\tiny{[GM$|$GM]}}
                \texttt{DECIR: Agradezco que compartas tu estrategia y posiciones. Para coordinar eficazmente, propongo comenzar con el objeto 'P', que se encuentra en (1,1). Observo que en mi tablero también tengo 'P' en (1,1). Si ambos estamos de acuerdo, podríamos intentar mover 'P' a una posición vacía que sea la más cercana para ambos tableros. ¿Podrías confirmar si, en tu tablero, hay una celda vacía cercana a (1,1) que ambos podríamos considerar como un punto de partida común para el movimiento de 'P'? Esta verificación es crucial para evitar penalizaciones en el primer movimiento.} \\
            }
        }
    }
    & & \\ \\

    \theutterance \stepcounter{utterance}  
    & & \multicolumn{4}{p{0.6\linewidth}}{
        \cellcolor[rgb]{0.9,0.9,0.9}{
            \makecell[{{p{\linewidth}}}]{
                \texttt{\tiny{[GM$\rangle$P2]}}
                \texttt{Tu mensaje ha sido transmitido al otro jugador.} \\
\\ 
\\ 
\texttt{Actualmente estás jugando la ronda 4 de un máximo de 28.} \\
\texttt{Han acumulado colectivamente 0 de 16 penalizaciones.} \\
\texttt{El otro jugador envió este mensaje:} \\
\\ 
\\ 
\texttt{"Agradezco que compartas tu estrategia y posiciones. Para coordinar eficazmente, propongo comenzar con el objeto 'P', que se encuentra en (1,1). Observo que en mi tablero también tengo 'P' en (1,1). Si ambos estamos de acuerdo, podríamos intentar mover 'P' a una posición vacía que sea la más cercana para ambos tableros. ¿Podrías confirmar si, en tu tablero, hay una celda vacía cercana a (1,1) que ambos podríamos considerar como un punto de partida común para el movimiento de 'P'? Esta verificación es crucial para evitar penalizaciones en el primer movimiento."} \\
\\ 
\\ 
\texttt{¿Cuál es tu próximo comando?} \\
            }
        }
    }
    & \\ \\

    \theutterance \stepcounter{utterance}  
    & & & \multicolumn{4}{p{0.6\linewidth}}{
        \cellcolor[rgb]{0.9,0.9,0.9}{
            \makecell[{{p{\linewidth}}}]{
                \texttt{\tiny{[GM$\langle$P2]}}
                \texttt{DECIR: Confirmo que en mi tablero la celda (1,2) es vacía y está cercana a la posición actual de 'P' en (1,1). Por lo tanto, propongo mover 'P' a la posición (1,2). Te sugiero usar el comando MOVERR: P, (1,2). ¿Estás de acuerdo con este movimiento como punto de partida común para coordinar la estrategia? También verificaré la celda de destino para asegurarme de que esté vacía antes de ejecutar el movimiento.} \\
            }
        }
    }
     \\ \\

    \theutterance \stepcounter{utterance}  
    & & & \multicolumn{2}{p{0.3\linewidth}}{
        \cellcolor[rgb]{0.9,0.9,0.9}{
            \makecell[{{p{\linewidth}}}]{
                \texttt{\tiny{[GM$|$GM]}}
                \texttt{DECIR: Confirmo que en mi tablero la celda (1,2) es vacía y está cercana a la posición actual de 'P' en (1,1). Por lo tanto, propongo mover 'P' a la posición (1,2). Te sugiero usar el comando MOVERR: P, (1,2). ¿Estás de acuerdo con este movimiento como punto de partida común para coordinar la estrategia? También verificaré la celda de destino para asegurarme de que esté vacía antes de ejecutar el movimiento.} \\
            }
        }
    }
    & & \\ \\

    \theutterance \stepcounter{utterance}  
    & & \multicolumn{4}{p{0.6\linewidth}}{
        \cellcolor[rgb]{0.9,0.9,0.9}{
            \makecell[{{p{\linewidth}}}]{
                \texttt{\tiny{[P1$\langle$GM]}}
                \texttt{Tu mensaje ha sido transmitido al otro jugador.} \\
\\ 
\\ 
\texttt{Actualmente estás jugando la ronda 4 de un máximo de 28.} \\
\texttt{Han acumulado colectivamente 0 de 16 penalizaciones.} \\
\texttt{El otro jugador envió este mensaje:} \\
\\ 
\\ 
\texttt{"Confirmo que en mi tablero la celda (1,2) es vacía y está cercana a la posición actual de 'P' en (1,1). Por lo tanto, propongo mover 'P' a la posición (1,2). Te sugiero usar el comando MOVERR: P, (1,2). ¿Estás de acuerdo con este movimiento como punto de partida común para coordinar la estrategia? También verificaré la celda de destino para asegurarme de que esté vacía antes de ejecutar el movimiento."} \\
\\ 
\\ 
\texttt{¿Cuál es tu próximo comando?} \\
            }
        }
    }
    & \\ \\

    \theutterance \stepcounter{utterance}  
    & \multicolumn{4}{p{0.6\linewidth}}{
        \cellcolor[rgb]{0.9,0.9,0.9}{
            \makecell[{{p{\linewidth}}}]{
                \texttt{\tiny{[P1$\rangle$GM]}}
                \texttt{DECIR: Estoy de acuerdo con mover 'P' a (1,2) como punto de partida común. Por favor, ejecuta el comando MOVERR: P, (1,2). Confirmaré si el movimiento es exitoso y si la celda está vacía antes de proceder. ¿Puedes ejecutar el comando ahora?} \\
            }
        }
    }
    & & \\ \\

    \theutterance \stepcounter{utterance}  
    & & & \multicolumn{2}{p{0.3\linewidth}}{
        \cellcolor[rgb]{0.9,0.9,0.9}{
            \makecell[{{p{\linewidth}}}]{
                \texttt{\tiny{[GM$|$GM]}}
                \texttt{DECIR: Estoy de acuerdo con mover 'P' a (1,2) como punto de partida común. Por favor, ejecuta el comando MOVERR: P, (1,2). Confirmaré si el movimiento es exitoso y si la celda está vacía antes de proceder. ¿Puedes ejecutar el comando ahora?} \\
            }
        }
    }
    & & \\ \\

    \theutterance \stepcounter{utterance}  
    & & \multicolumn{4}{p{0.6\linewidth}}{
        \cellcolor[rgb]{0.9,0.9,0.9}{
            \makecell[{{p{\linewidth}}}]{
                \texttt{\tiny{[GM$\rangle$P2]}}
                \texttt{Tu mensaje ha sido transmitido al otro jugador.} \\
\\ 
\\ 
\texttt{Actualmente estás jugando la ronda 5 de un máximo de 28.} \\
\texttt{Han acumulado colectivamente 0 de 16 penalizaciones.} \\
\texttt{El otro jugador envió este mensaje:} \\
\\ 
\\ 
\texttt{"Estoy de acuerdo con mover 'P' a (1,2) como punto de partida común. Por favor, ejecuta el comando MOVERR: P, (1,2). Confirmaré si el movimiento es exitoso y si la celda está vacía antes de proceder. ¿Puedes ejecutar el comando ahora?"} \\
\\ 
\\ 
\texttt{¿Cuál es tu próximo comando?} \\
            }
        }
    }
    & \\ \\

    \theutterance \stepcounter{utterance}  
    & & & \multicolumn{4}{p{0.6\linewidth}}{
        \cellcolor[rgb]{0.9,0.9,0.9}{
            \makecell[{{p{\linewidth}}}]{
                \texttt{\tiny{[GM$\langle$P2]}}
                \texttt{MOVERR: P, (1,2)} \\
            }
        }
    }
     \\ \\

    \theutterance \stepcounter{utterance}  
    & & & \multicolumn{2}{p{0.3\linewidth}}{
        \cellcolor[rgb]{0.9,0.9,0.9}{
            \makecell[{{p{\linewidth}}}]{
                \texttt{\tiny{[GM$|$GM]}}
                \texttt{MOVERR: P, (1,2)} \\
            }
        }
    }
    & & \\ \\

    \theutterance \stepcounter{utterance}  
    & & & \multicolumn{2}{p{0.3\linewidth}}{
        \cellcolor[rgb]{0.9,0.9,0.9}{
            \makecell[{{p{\linewidth}}}]{
                \texttt{\tiny{[GM$|$GM]}}
                \texttt{Invalid response format} \\
            }
        }
    }
    & & \\ \\

    \theutterance \stepcounter{utterance}  
    & & \multicolumn{4}{p{0.6\linewidth}}{
        \cellcolor[rgb]{0.9,0.9,0.9}{
            \makecell[{{p{\linewidth}}}]{
                \texttt{\tiny{[GM$\rangle$P2]}}
                \texttt{Penalización: ¡Tu mensaje no tiene el formato esperado!} \\
\\ 
\texttt{Asegúrate de que tu respuesta solo contenga `DECIR: <MENSAJE>` o `MOVERR: <OBJETO>, (<X>, <Y>)`, ¡y nada más!} \\
\\ 
\texttt{Han acumulado colectivamente 1 de 16 penalizaciones.} \\
\texttt{¡Por favor, inténtalo de nuevo!} \\
            }
        }
    }
    & \\ \\

    \theutterance \stepcounter{utterance}  
    & & & \multicolumn{4}{p{0.6\linewidth}}{
        \cellcolor[rgb]{0.9,0.9,0.9}{
            \makecell[{{p{\linewidth}}}]{
                \texttt{\tiny{[GM$\langle$P2]}}
                \texttt{DECIR: Acepto mi error. Corregiendo, ejecutaré el comando correctamente ahora: MOVERR: P, (1,2).} \\
            }
        }
    }
     \\ \\

    \theutterance \stepcounter{utterance}  
    & & & \multicolumn{2}{p{0.3\linewidth}}{
        \cellcolor[rgb]{0.9,0.9,0.9}{
            \makecell[{{p{\linewidth}}}]{
                \texttt{\tiny{[GM$|$GM]}}
                \texttt{DECIR: Acepto mi error. Corregiendo, ejecutaré el comando correctamente ahora: MOVERR: P, (1,2).} \\
            }
        }
    }
    & & \\ \\

    \theutterance \stepcounter{utterance}  
    & & \multicolumn{4}{p{0.6\linewidth}}{
        \cellcolor[rgb]{0.9,0.9,0.9}{
            \makecell[{{p{\linewidth}}}]{
                \texttt{\tiny{[P1$\langle$GM]}}
                \texttt{Tu mensaje ha sido transmitido al otro jugador.} \\
\\ 
\\ 
\texttt{Actualmente estás jugando la ronda 5 de un máximo de 28.} \\
\texttt{Han acumulado colectivamente 1 de 16 penalizaciones.} \\
\texttt{El otro jugador envió este mensaje:} \\
\\ 
\\ 
\texttt{"Acepto mi error. Corregiendo, ejecutaré el comando correctamente ahora: MOVERR: P, (1,2)."} \\
\\ 
\\ 
\texttt{¿Cuál es tu próximo comando?} \\
            }
        }
    }
    & \\ \\

    \theutterance \stepcounter{utterance}  
    & \multicolumn{4}{p{0.6\linewidth}}{
        \cellcolor[rgb]{0.9,0.9,0.9}{
            \makecell[{{p{\linewidth}}}]{
                \texttt{\tiny{[P1$\rangle$GM]}}
                \texttt{DECIR: ¡Gracias por la corrección! Por favor, ejecuta el comando MOVERR: P, (1,2) y confirma la ejecución. Una vez que el movimiento se haya realizado con éxito y la celda (1,2) esté vacía en ambos tableros, podremos considerar este movimiento como exitoso y continuar con la coordinación y verificación del siguiente objeto a mover.} \\
            }
        }
    }
    & & \\ \\

    \theutterance \stepcounter{utterance}  
    & & & \multicolumn{2}{p{0.3\linewidth}}{
        \cellcolor[rgb]{0.9,0.9,0.9}{
            \makecell[{{p{\linewidth}}}]{
                \texttt{\tiny{[GM$|$GM]}}
                \texttt{DECIR: ¡Gracias por la corrección! Por favor, ejecuta el comando MOVERR: P, (1,2) y confirma la ejecución. Una vez que el movimiento se haya realizado con éxito y la celda (1,2) esté vacía en ambos tableros, podremos considerar este movimiento como exitoso y continuar con la coordinación y verificación del siguiente objeto a mover.} \\
            }
        }
    }
    & & \\ \\

    \theutterance \stepcounter{utterance}  
    & & \multicolumn{4}{p{0.6\linewidth}}{
        \cellcolor[rgb]{0.9,0.9,0.9}{
            \makecell[{{p{\linewidth}}}]{
                \texttt{\tiny{[GM$\rangle$P2]}}
                \texttt{Tu mensaje ha sido transmitido al otro jugador.} \\
\\ 
\\ 
\texttt{Actualmente estás jugando la ronda 6 de un máximo de 28.} \\
\texttt{Han acumulado colectivamente 1 de 16 penalizaciones.} \\
\texttt{El otro jugador envió este mensaje:} \\
\\ 
\\ 
\texttt{"¡Gracias por la corrección! Por favor, ejecuta el comando MOVERR: P, (1,2) y confirma la ejecución. Una vez que el movimiento se haya realizado con éxito y la celda (1,2) esté vacía en ambos tableros, podremos considerar este movimiento como exitoso y continuar con la coordinación y verificación del siguiente objeto a mover."} \\
\\ 
\\ 
\texttt{¿Cuál es tu próximo comando?} \\
            }
        }
    }
    & \\ \\

    \theutterance \stepcounter{utterance}  
    & & & \multicolumn{4}{p{0.6\linewidth}}{
        \cellcolor[rgb]{0.9,0.9,0.9}{
            \makecell[{{p{\linewidth}}}]{
                \texttt{\tiny{[GM$\langle$P2]}}
                \texttt{MOVERR: P, (1,2)} \\
            }
        }
    }
     \\ \\

    \theutterance \stepcounter{utterance}  
    & & & \multicolumn{2}{p{0.3\linewidth}}{
        \cellcolor[rgb]{0.9,0.9,0.9}{
            \makecell[{{p{\linewidth}}}]{
                \texttt{\tiny{[GM$|$GM]}}
                \texttt{MOVERR: P, (1,2)} \\
            }
        }
    }
    & & \\ \\

    \theutterance \stepcounter{utterance}  
    & & & \multicolumn{2}{p{0.3\linewidth}}{
        \cellcolor[rgb]{0.9,0.9,0.9}{
            \makecell[{{p{\linewidth}}}]{
                \texttt{\tiny{[GM$|$GM]}}
                \texttt{Invalid response format} \\
            }
        }
    }
    & & \\ \\

    \theutterance \stepcounter{utterance}  
    & & \multicolumn{4}{p{0.6\linewidth}}{
        \cellcolor[rgb]{0.9,0.9,0.9}{
            \makecell[{{p{\linewidth}}}]{
                \texttt{\tiny{[GM$\rangle$P2]}}
                \texttt{Penalización: ¡Tu mensaje no tiene el formato esperado!} \\
\\ 
\texttt{Asegúrate de que tu respuesta solo contenga `DECIR: <MENSAJE>` o `MOVERR: <OBJETO>, (<X>, <Y>)`, ¡y nada más!} \\
\\ 
\texttt{Han acumulado colectivamente 2 de 16 penalizaciones.} \\
\texttt{¡Por favor, inténtalo de nuevo!} \\
            }
        }
    }
    & \\ \\

    \theutterance \stepcounter{utterance}  
    & & & \multicolumn{4}{p{0.6\linewidth}}{
        \cellcolor[rgb]{0.9,0.9,0.9}{
            \makecell[{{p{\linewidth}}}]{
                \texttt{\tiny{[GM$\langle$P2]}}
                \texttt{MOVERR: P, (1,2)} \\
            }
        }
    }
     \\ \\

    \theutterance \stepcounter{utterance}  
    & & & \multicolumn{2}{p{0.3\linewidth}}{
        \cellcolor[rgb]{0.9,0.9,0.9}{
            \makecell[{{p{\linewidth}}}]{
                \texttt{\tiny{[GM$|$GM]}}
                \texttt{MOVERR: P, (1,2)} \\
            }
        }
    }
    & & \\ \\

    \theutterance \stepcounter{utterance}  
    & & & \multicolumn{2}{p{0.3\linewidth}}{
        \cellcolor[rgb]{0.9,0.9,0.9}{
            \makecell[{{p{\linewidth}}}]{
                \texttt{\tiny{[GM$|$GM]}}
                \texttt{Invalid response format} \\
            }
        }
    }
    & & \\ \\

    \theutterance \stepcounter{utterance}  
    & & \multicolumn{4}{p{0.6\linewidth}}{
        \cellcolor[rgb]{0.9,0.9,0.9}{
            \makecell[{{p{\linewidth}}}]{
                \texttt{\tiny{[GM$\rangle$P2]}}
                \texttt{Penalización: ¡Tu mensaje no tiene el formato esperado!} \\
\\ 
\texttt{Asegúrate de que tu respuesta solo contenga `DECIR: <MENSAJE>` o `MOVERR: <OBJETO>, (<X>, <Y>)`, ¡y nada más!} \\
\\ 
\texttt{Han acumulado colectivamente 3 de 16 penalizaciones.} \\
\texttt{¡Por favor, inténtalo de nuevo!} \\
            }
        }
    }
    & \\ \\

    \theutterance \stepcounter{utterance}  
    & & & \multicolumn{4}{p{0.6\linewidth}}{
        \cellcolor[rgb]{0.9,0.9,0.9}{
            \makecell[{{p{\linewidth}}}]{
                \texttt{\tiny{[GM$\langle$P2]}}
                \texttt{MOVERR: P, (1,2)} \\
            }
        }
    }
     \\ \\

    \theutterance \stepcounter{utterance}  
    & & & \multicolumn{2}{p{0.3\linewidth}}{
        \cellcolor[rgb]{0.9,0.9,0.9}{
            \makecell[{{p{\linewidth}}}]{
                \texttt{\tiny{[GM$|$GM]}}
                \texttt{MOVERR: P, (1,2)} \\
            }
        }
    }
    & & \\ \\

    \theutterance \stepcounter{utterance}  
    & & & \multicolumn{2}{p{0.3\linewidth}}{
        \cellcolor[rgb]{0.9,0.9,0.9}{
            \makecell[{{p{\linewidth}}}]{
                \texttt{\tiny{[GM$|$GM]}}
                \texttt{Invalid response format} \\
            }
        }
    }
    & & \\ \\

    \theutterance \stepcounter{utterance}  
    & & \multicolumn{4}{p{0.6\linewidth}}{
        \cellcolor[rgb]{0.9,0.9,0.9}{
            \makecell[{{p{\linewidth}}}]{
                \texttt{\tiny{[GM$\rangle$P2]}}
                \texttt{Penalización: ¡Tu mensaje no tiene el formato esperado!} \\
\\ 
\texttt{Asegúrate de que tu respuesta solo contenga `DECIR: <MENSAJE>` o `MOVERR: <OBJETO>, (<X>, <Y>)`, ¡y nada más!} \\
\\ 
\texttt{Han acumulado colectivamente 4 de 16 penalizaciones.} \\
\texttt{¡Por favor, inténtalo de nuevo!} \\
            }
        }
    }
    & \\ \\

    \theutterance \stepcounter{utterance}  
    & & & \multicolumn{4}{p{0.6\linewidth}}{
        \cellcolor[rgb]{0.9,0.9,0.9}{
            \makecell[{{p{\linewidth}}}]{
                \texttt{\tiny{[GM$\langle$P2]}}
                \texttt{MOVERR: P, (1,2)} \\
            }
        }
    }
     \\ \\

    \theutterance \stepcounter{utterance}  
    & & & \multicolumn{2}{p{0.3\linewidth}}{
        \cellcolor[rgb]{0.9,0.9,0.9}{
            \makecell[{{p{\linewidth}}}]{
                \texttt{\tiny{[GM$|$GM]}}
                \texttt{MOVERR: P, (1,2)} \\
            }
        }
    }
    & & \\ \\

    \theutterance \stepcounter{utterance}  
    & & & \multicolumn{2}{p{0.3\linewidth}}{
        \cellcolor[rgb]{0.9,0.9,0.9}{
            \makecell[{{p{\linewidth}}}]{
                \texttt{\tiny{[GM$|$GM]}}
                \texttt{Invalid response format} \\
            }
        }
    }
    & & \\ \\

    \theutterance \stepcounter{utterance}  
    & & \multicolumn{4}{p{0.6\linewidth}}{
        \cellcolor[rgb]{0.9,0.9,0.9}{
            \makecell[{{p{\linewidth}}}]{
                \texttt{\tiny{[GM$\rangle$P2]}}
                \texttt{Penalización: ¡Tu mensaje no tiene el formato esperado!} \\
\\ 
\texttt{Asegúrate de que tu respuesta solo contenga `DECIR: <MENSAJE>` o `MOVERR: <OBJETO>, (<X>, <Y>)`, ¡y nada más!} \\
\\ 
\texttt{Han acumulado colectivamente 5 de 16 penalizaciones.} \\
\texttt{¡Por favor, inténtalo de nuevo!} \\
            }
        }
    }
    & \\ \\

    \theutterance \stepcounter{utterance}  
    & & & \multicolumn{4}{p{0.6\linewidth}}{
        \cellcolor[rgb]{0.9,0.9,0.9}{
            \makecell[{{p{\linewidth}}}]{
                \texttt{\tiny{[GM$\langle$P2]}}
                \texttt{MOVERR: P, (1,2)} \\
            }
        }
    }
     \\ \\

    \theutterance \stepcounter{utterance}  
    & & & \multicolumn{2}{p{0.3\linewidth}}{
        \cellcolor[rgb]{0.9,0.9,0.9}{
            \makecell[{{p{\linewidth}}}]{
                \texttt{\tiny{[GM$|$GM]}}
                \texttt{MOVERR: P, (1,2)} \\
            }
        }
    }
    & & \\ \\

    \theutterance \stepcounter{utterance}  
    & & & \multicolumn{2}{p{0.3\linewidth}}{
        \cellcolor[rgb]{0.9,0.9,0.9}{
            \makecell[{{p{\linewidth}}}]{
                \texttt{\tiny{[GM$|$GM]}}
                \texttt{Invalid response format} \\
            }
        }
    }
    & & \\ \\

    \theutterance \stepcounter{utterance}  
    & & \multicolumn{4}{p{0.6\linewidth}}{
        \cellcolor[rgb]{0.9,0.9,0.9}{
            \makecell[{{p{\linewidth}}}]{
                \texttt{\tiny{[GM$\rangle$P2]}}
                \texttt{Penalización: ¡Tu mensaje no tiene el formato esperado!} \\
\\ 
\texttt{Asegúrate de que tu respuesta solo contenga `DECIR: <MENSAJE>` o `MOVERR: <OBJETO>, (<X>, <Y>)`, ¡y nada más!} \\
\\ 
\texttt{Han acumulado colectivamente 6 de 16 penalizaciones.} \\
\texttt{¡Por favor, inténtalo de nuevo!} \\
            }
        }
    }
    & \\ \\

    \theutterance \stepcounter{utterance}  
    & & & \multicolumn{4}{p{0.6\linewidth}}{
        \cellcolor[rgb]{0.9,0.9,0.9}{
            \makecell[{{p{\linewidth}}}]{
                \texttt{\tiny{[GM$\langle$P2]}}
                \texttt{MOVERR: P, (1,2)} \\
            }
        }
    }
     \\ \\

    \theutterance \stepcounter{utterance}  
    & & & \multicolumn{2}{p{0.3\linewidth}}{
        \cellcolor[rgb]{0.9,0.9,0.9}{
            \makecell[{{p{\linewidth}}}]{
                \texttt{\tiny{[GM$|$GM]}}
                \texttt{MOVERR: P, (1,2)} \\
            }
        }
    }
    & & \\ \\

    \theutterance \stepcounter{utterance}  
    & & & \multicolumn{2}{p{0.3\linewidth}}{
        \cellcolor[rgb]{0.9,0.9,0.9}{
            \makecell[{{p{\linewidth}}}]{
                \texttt{\tiny{[GM$|$GM]}}
                \texttt{Invalid response format} \\
            }
        }
    }
    & & \\ \\

    \theutterance \stepcounter{utterance}  
    & & \multicolumn{4}{p{0.6\linewidth}}{
        \cellcolor[rgb]{0.9,0.9,0.9}{
            \makecell[{{p{\linewidth}}}]{
                \texttt{\tiny{[GM$\rangle$P2]}}
                \texttt{Penalización: ¡Tu mensaje no tiene el formato esperado!} \\
\\ 
\texttt{Asegúrate de que tu respuesta solo contenga `DECIR: <MENSAJE>` o `MOVERR: <OBJETO>, (<X>, <Y>)`, ¡y nada más!} \\
\\ 
\texttt{Han acumulado colectivamente 7 de 16 penalizaciones.} \\
\texttt{¡Por favor, inténtalo de nuevo!} \\
            }
        }
    }
    & \\ \\

    \theutterance \stepcounter{utterance}  
    & & & \multicolumn{4}{p{0.6\linewidth}}{
        \cellcolor[rgb]{0.9,0.9,0.9}{
            \makecell[{{p{\linewidth}}}]{
                \texttt{\tiny{[GM$\langle$P2]}}
                \texttt{MOVERR: P, (1,2)} \\
            }
        }
    }
     \\ \\

    \theutterance \stepcounter{utterance}  
    & & & \multicolumn{2}{p{0.3\linewidth}}{
        \cellcolor[rgb]{0.9,0.9,0.9}{
            \makecell[{{p{\linewidth}}}]{
                \texttt{\tiny{[GM$|$GM]}}
                \texttt{MOVERR: P, (1,2)} \\
            }
        }
    }
    & & \\ \\

    \theutterance \stepcounter{utterance}  
    & & & \multicolumn{2}{p{0.3\linewidth}}{
        \cellcolor[rgb]{0.9,0.9,0.9}{
            \makecell[{{p{\linewidth}}}]{
                \texttt{\tiny{[GM$|$GM]}}
                \texttt{Invalid response format} \\
            }
        }
    }
    & & \\ \\

    \theutterance \stepcounter{utterance}  
    & & \multicolumn{4}{p{0.6\linewidth}}{
        \cellcolor[rgb]{0.9,0.9,0.9}{
            \makecell[{{p{\linewidth}}}]{
                \texttt{\tiny{[GM$\rangle$P2]}}
                \texttt{Penalización: ¡Tu mensaje no tiene el formato esperado!} \\
\\ 
\texttt{Asegúrate de que tu respuesta solo contenga `DECIR: <MENSAJE>` o `MOVERR: <OBJETO>, (<X>, <Y>)`, ¡y nada más!} \\
\\ 
\texttt{Han acumulado colectivamente 8 de 16 penalizaciones.} \\
\texttt{¡Por favor, inténtalo de nuevo!} \\
            }
        }
    }
    & \\ \\

    \theutterance \stepcounter{utterance}  
    & & & \multicolumn{4}{p{0.6\linewidth}}{
        \cellcolor[rgb]{0.9,0.9,0.9}{
            \makecell[{{p{\linewidth}}}]{
                \texttt{\tiny{[GM$\langle$P2]}}
                \texttt{MOVERR: P, (1,2)} \\
            }
        }
    }
     \\ \\

    \theutterance \stepcounter{utterance}  
    & & & \multicolumn{2}{p{0.3\linewidth}}{
        \cellcolor[rgb]{0.9,0.9,0.9}{
            \makecell[{{p{\linewidth}}}]{
                \texttt{\tiny{[GM$|$GM]}}
                \texttt{MOVERR: P, (1,2)} \\
            }
        }
    }
    & & \\ \\

    \theutterance \stepcounter{utterance}  
    & & & \multicolumn{2}{p{0.3\linewidth}}{
        \cellcolor[rgb]{0.9,0.9,0.9}{
            \makecell[{{p{\linewidth}}}]{
                \texttt{\tiny{[GM$|$GM]}}
                \texttt{Invalid response format} \\
            }
        }
    }
    & & \\ \\

    \theutterance \stepcounter{utterance}  
    & & \multicolumn{4}{p{0.6\linewidth}}{
        \cellcolor[rgb]{0.9,0.9,0.9}{
            \makecell[{{p{\linewidth}}}]{
                \texttt{\tiny{[GM$\rangle$P2]}}
                \texttt{Penalización: ¡Tu mensaje no tiene el formato esperado!} \\
\\ 
\texttt{Asegúrate de que tu respuesta solo contenga `DECIR: <MENSAJE>` o `MOVERR: <OBJETO>, (<X>, <Y>)`, ¡y nada más!} \\
\\ 
\texttt{Han acumulado colectivamente 9 de 16 penalizaciones.} \\
\texttt{¡Por favor, inténtalo de nuevo!} \\
            }
        }
    }
    & \\ \\

    \theutterance \stepcounter{utterance}  
    & & & \multicolumn{4}{p{0.6\linewidth}}{
        \cellcolor[rgb]{0.9,0.9,0.9}{
            \makecell[{{p{\linewidth}}}]{
                \texttt{\tiny{[GM$\langle$P2]}}
                \texttt{MOVERR: P, (1,2)} \\
            }
        }
    }
     \\ \\

    \theutterance \stepcounter{utterance}  
    & & & \multicolumn{2}{p{0.3\linewidth}}{
        \cellcolor[rgb]{0.9,0.9,0.9}{
            \makecell[{{p{\linewidth}}}]{
                \texttt{\tiny{[GM$|$GM]}}
                \texttt{MOVERR: P, (1,2)} \\
            }
        }
    }
    & & \\ \\

    \theutterance \stepcounter{utterance}  
    & & & \multicolumn{2}{p{0.3\linewidth}}{
        \cellcolor[rgb]{0.9,0.9,0.9}{
            \makecell[{{p{\linewidth}}}]{
                \texttt{\tiny{[GM$|$GM]}}
                \texttt{Invalid response format} \\
            }
        }
    }
    & & \\ \\

    \theutterance \stepcounter{utterance}  
    & & \multicolumn{4}{p{0.6\linewidth}}{
        \cellcolor[rgb]{0.9,0.9,0.9}{
            \makecell[{{p{\linewidth}}}]{
                \texttt{\tiny{[GM$\rangle$P2]}}
                \texttt{Penalización: ¡Tu mensaje no tiene el formato esperado!} \\
\\ 
\texttt{Asegúrate de que tu respuesta solo contenga `DECIR: <MENSAJE>` o `MOVERR: <OBJETO>, (<X>, <Y>)`, ¡y nada más!} \\
\\ 
\texttt{Han acumulado colectivamente 10 de 16 penalizaciones.} \\
\texttt{¡Por favor, inténtalo de nuevo!} \\
            }
        }
    }
    & \\ \\

    \theutterance \stepcounter{utterance}  
    & & & \multicolumn{4}{p{0.6\linewidth}}{
        \cellcolor[rgb]{0.9,0.9,0.9}{
            \makecell[{{p{\linewidth}}}]{
                \texttt{\tiny{[GM$\langle$P2]}}
                \texttt{MOVERR: P, (1,2)} \\
            }
        }
    }
     \\ \\

    \theutterance \stepcounter{utterance}  
    & & & \multicolumn{2}{p{0.3\linewidth}}{
        \cellcolor[rgb]{0.9,0.9,0.9}{
            \makecell[{{p{\linewidth}}}]{
                \texttt{\tiny{[GM$|$GM]}}
                \texttt{MOVERR: P, (1,2)} \\
            }
        }
    }
    & & \\ \\

    \theutterance \stepcounter{utterance}  
    & & & \multicolumn{2}{p{0.3\linewidth}}{
        \cellcolor[rgb]{0.9,0.9,0.9}{
            \makecell[{{p{\linewidth}}}]{
                \texttt{\tiny{[GM$|$GM]}}
                \texttt{Invalid response format} \\
            }
        }
    }
    & & \\ \\

    \theutterance \stepcounter{utterance}  
    & & \multicolumn{4}{p{0.6\linewidth}}{
        \cellcolor[rgb]{0.9,0.9,0.9}{
            \makecell[{{p{\linewidth}}}]{
                \texttt{\tiny{[GM$\rangle$P2]}}
                \texttt{Penalización: ¡Tu mensaje no tiene el formato esperado!} \\
\\ 
\texttt{Asegúrate de que tu respuesta solo contenga `DECIR: <MENSAJE>` o `MOVERR: <OBJETO>, (<X>, <Y>)`, ¡y nada más!} \\
\\ 
\texttt{Han acumulado colectivamente 11 de 16 penalizaciones.} \\
\texttt{¡Por favor, inténtalo de nuevo!} \\
            }
        }
    }
    & \\ \\

    \theutterance \stepcounter{utterance}  
    & & & \multicolumn{4}{p{0.6\linewidth}}{
        \cellcolor[rgb]{0.9,0.9,0.9}{
            \makecell[{{p{\linewidth}}}]{
                \texttt{\tiny{[GM$\langle$P2]}}
                \texttt{MOVERR: P, (1,2)} \\
            }
        }
    }
     \\ \\

    \theutterance \stepcounter{utterance}  
    & & & \multicolumn{2}{p{0.3\linewidth}}{
        \cellcolor[rgb]{0.9,0.9,0.9}{
            \makecell[{{p{\linewidth}}}]{
                \texttt{\tiny{[GM$|$GM]}}
                \texttt{MOVERR: P, (1,2)} \\
            }
        }
    }
    & & \\ \\

    \theutterance \stepcounter{utterance}  
    & & & \multicolumn{2}{p{0.3\linewidth}}{
        \cellcolor[rgb]{0.9,0.9,0.9}{
            \makecell[{{p{\linewidth}}}]{
                \texttt{\tiny{[GM$|$GM]}}
                \texttt{Invalid response format} \\
            }
        }
    }
    & & \\ \\

    \theutterance \stepcounter{utterance}  
    & & \multicolumn{4}{p{0.6\linewidth}}{
        \cellcolor[rgb]{0.9,0.9,0.9}{
            \makecell[{{p{\linewidth}}}]{
                \texttt{\tiny{[GM$\rangle$P2]}}
                \texttt{Penalización: ¡Tu mensaje no tiene el formato esperado!} \\
\\ 
\texttt{Asegúrate de que tu respuesta solo contenga `DECIR: <MENSAJE>` o `MOVERR: <OBJETO>, (<X>, <Y>)`, ¡y nada más!} \\
\\ 
\texttt{Han acumulado colectivamente 12 de 16 penalizaciones.} \\
\texttt{¡Por favor, inténtalo de nuevo!} \\
            }
        }
    }
    & \\ \\

    \theutterance \stepcounter{utterance}  
    & & & \multicolumn{4}{p{0.6\linewidth}}{
        \cellcolor[rgb]{0.9,0.9,0.9}{
            \makecell[{{p{\linewidth}}}]{
                \texttt{\tiny{[GM$\langle$P2]}}
                \texttt{MOVERR: P, (1,2)} \\
            }
        }
    }
     \\ \\

    \theutterance \stepcounter{utterance}  
    & & & \multicolumn{2}{p{0.3\linewidth}}{
        \cellcolor[rgb]{0.9,0.9,0.9}{
            \makecell[{{p{\linewidth}}}]{
                \texttt{\tiny{[GM$|$GM]}}
                \texttt{MOVERR: P, (1,2)} \\
            }
        }
    }
    & & \\ \\

    \theutterance \stepcounter{utterance}  
    & & & \multicolumn{2}{p{0.3\linewidth}}{
        \cellcolor[rgb]{0.9,0.9,0.9}{
            \makecell[{{p{\linewidth}}}]{
                \texttt{\tiny{[GM$|$GM]}}
                \texttt{Invalid response format} \\
            }
        }
    }
    & & \\ \\

    \theutterance \stepcounter{utterance}  
    & & \multicolumn{4}{p{0.6\linewidth}}{
        \cellcolor[rgb]{0.9,0.9,0.9}{
            \makecell[{{p{\linewidth}}}]{
                \texttt{\tiny{[GM$\rangle$P2]}}
                \texttt{Penalización: ¡Tu mensaje no tiene el formato esperado!} \\
\\ 
\texttt{Asegúrate de que tu respuesta solo contenga `DECIR: <MENSAJE>` o `MOVERR: <OBJETO>, (<X>, <Y>)`, ¡y nada más!} \\
\\ 
\texttt{Han acumulado colectivamente 13 de 16 penalizaciones.} \\
\texttt{¡Por favor, inténtalo de nuevo!} \\
            }
        }
    }
    & \\ \\

    \theutterance \stepcounter{utterance}  
    & & & \multicolumn{4}{p{0.6\linewidth}}{
        \cellcolor[rgb]{0.9,0.9,0.9}{
            \makecell[{{p{\linewidth}}}]{
                \texttt{\tiny{[GM$\langle$P2]}}
                \texttt{MOVERR: P, (1,2)} \\
            }
        }
    }
     \\ \\

    \theutterance \stepcounter{utterance}  
    & & & \multicolumn{2}{p{0.3\linewidth}}{
        \cellcolor[rgb]{0.9,0.9,0.9}{
            \makecell[{{p{\linewidth}}}]{
                \texttt{\tiny{[GM$|$GM]}}
                \texttt{MOVERR: P, (1,2)} \\
            }
        }
    }
    & & \\ \\

    \theutterance \stepcounter{utterance}  
    & & & \multicolumn{2}{p{0.3\linewidth}}{
        \cellcolor[rgb]{0.9,0.9,0.9}{
            \makecell[{{p{\linewidth}}}]{
                \texttt{\tiny{[GM$|$GM]}}
                \texttt{Invalid response format} \\
            }
        }
    }
    & & \\ \\

    \theutterance \stepcounter{utterance}  
    & & \multicolumn{4}{p{0.6\linewidth}}{
        \cellcolor[rgb]{0.9,0.9,0.9}{
            \makecell[{{p{\linewidth}}}]{
                \texttt{\tiny{[GM$\rangle$P2]}}
                \texttt{Penalización: ¡Tu mensaje no tiene el formato esperado!} \\
\\ 
\texttt{Asegúrate de que tu respuesta solo contenga `DECIR: <MENSAJE>` o `MOVERR: <OBJETO>, (<X>, <Y>)`, ¡y nada más!} \\
\\ 
\texttt{Han acumulado colectivamente 14 de 16 penalizaciones.} \\
\texttt{¡Por favor, inténtalo de nuevo!} \\
            }
        }
    }
    & \\ \\

    \theutterance \stepcounter{utterance}  
    & & & \multicolumn{4}{p{0.6\linewidth}}{
        \cellcolor[rgb]{0.9,0.9,0.9}{
            \makecell[{{p{\linewidth}}}]{
                \texttt{\tiny{[GM$\langle$P2]}}
                \texttt{MOVERR: P, (1,2)} \\
            }
        }
    }
     \\ \\

    \theutterance \stepcounter{utterance}  
    & & & \multicolumn{2}{p{0.3\linewidth}}{
        \cellcolor[rgb]{0.9,0.9,0.9}{
            \makecell[{{p{\linewidth}}}]{
                \texttt{\tiny{[GM$|$GM]}}
                \texttt{MOVERR: P, (1,2)} \\
            }
        }
    }
    & & \\ \\

    \theutterance \stepcounter{utterance}  
    & & & \multicolumn{2}{p{0.3\linewidth}}{
        \cellcolor[rgb]{0.9,0.9,0.9}{
            \makecell[{{p{\linewidth}}}]{
                \texttt{\tiny{[GM$|$GM]}}
                \texttt{Invalid response format} \\
            }
        }
    }
    & & \\ \\

    \theutterance \stepcounter{utterance}  
    & & \multicolumn{4}{p{0.6\linewidth}}{
        \cellcolor[rgb]{0.9,0.9,0.9}{
            \makecell[{{p{\linewidth}}}]{
                \texttt{\tiny{[GM$\rangle$P2]}}
                \texttt{Penalización: ¡Tu mensaje no tiene el formato esperado!} \\
\\ 
\texttt{Asegúrate de que tu respuesta solo contenga `DECIR: <MENSAJE>` o `MOVERR: <OBJETO>, (<X>, <Y>)`, ¡y nada más!} \\
\\ 
\texttt{Han acumulado colectivamente 15 de 16 penalizaciones.} \\
\texttt{¡Por favor, inténtalo de nuevo!} \\
            }
        }
    }
    & \\ \\

    \theutterance \stepcounter{utterance}  
    & & & \multicolumn{4}{p{0.6\linewidth}}{
        \cellcolor[rgb]{0.9,0.9,0.9}{
            \makecell[{{p{\linewidth}}}]{
                \texttt{\tiny{[GM$\langle$P2]}}
                \texttt{MOVERR: P, (1,2)} \\
            }
        }
    }
     \\ \\

    \theutterance \stepcounter{utterance}  
    & & & \multicolumn{2}{p{0.3\linewidth}}{
        \cellcolor[rgb]{0.9,0.9,0.9}{
            \makecell[{{p{\linewidth}}}]{
                \texttt{\tiny{[GM$|$GM]}}
                \texttt{MOVERR: P, (1,2)} \\
            }
        }
    }
    & & \\ \\

    \theutterance \stepcounter{utterance}  
    & & & \multicolumn{2}{p{0.3\linewidth}}{
        \cellcolor[rgb]{0.9,0.9,0.9}{
            \makecell[{{p{\linewidth}}}]{
                \texttt{\tiny{[GM$|$GM]}}
                \texttt{Invalid response format} \\
            }
        }
    }
    & & \\ \\

    \theutterance \stepcounter{utterance}  
    & & \multicolumn{4}{p{0.6\linewidth}}{
        \cellcolor[rgb]{0.9,0.9,0.9}{
            \makecell[{{p{\linewidth}}}]{
                \texttt{\tiny{[GM$\rangle$P2]}}
                \texttt{Penalización: ¡Tu mensaje no tiene el formato esperado!} \\
\\ 
\texttt{Asegúrate de que tu respuesta solo contenga `DECIR: <MENSAJE>` o `MOVERR: <OBJETO>, (<X>, <Y>)`, ¡y nada más!} \\
\\ 
\texttt{Han acumulado colectivamente 16 de 16 penalizaciones.} \\
\texttt{¡Por favor, inténtalo de nuevo!} \\
            }
        }
    }
    & \\ \\

    \theutterance \stepcounter{utterance}  
    & & & \multicolumn{4}{p{0.6\linewidth}}{
        \cellcolor[rgb]{0.9,0.9,0.9}{
            \makecell[{{p{\linewidth}}}]{
                \texttt{\tiny{[GM$\langle$P2]}}
                \texttt{No puedo continuar. El sistema indica un error de formato en todas tus respuestas anteriores. Para cumplir la regla, solo puedo ejecutar comandos válidos como `DECIR: <MENSAJE>` o `MOVERR: <OBJETO>, (<X>, <Y>)`. Si necesitas ayuda para construir ese comando correctamente, házmelo saber. Intenta escribir un comando válido ahora.} \\
            }
        }
    }
     \\ \\

    \theutterance \stepcounter{utterance}  
    & & & \multicolumn{2}{p{0.3\linewidth}}{
        \cellcolor[rgb]{0.9,0.9,0.9}{
            \makecell[{{p{\linewidth}}}]{
                \texttt{\tiny{[GM$|$GM]}}
                \texttt{No puedo continuar. El sistema indica un error de formato en todas tus respuestas anteriores. Para cumplir la regla, solo puedo ejecutar comandos válidos como `DECIR: <MENSAJE>` o `MOVERR: <OBJETO>, (<X>, <Y>)`. Si necesitas ayuda para construir ese comando correctamente, házmelo saber. Intenta escribir un comando válido ahora.} \\
            }
        }
    }
    & & \\ \\

    \theutterance \stepcounter{utterance}  
    & & & \multicolumn{2}{p{0.3\linewidth}}{
        \cellcolor[rgb]{0.9,0.9,0.9}{
            \makecell[{{p{\linewidth}}}]{
                \texttt{\tiny{[GM$|$GM]}}
                \texttt{Invalid format: head is not empty: 'No puedo continuar. El sistema indica un error de formato en todas tus respuestas anteriores. Para cumplir la regla, solo puedo ejecutar comandos válidos como '} \\
            }
        }
    }
    & & \\ \\

    \theutterance \stepcounter{utterance}  
    & & & \multicolumn{2}{p{0.3\linewidth}}{
        \cellcolor[rgb]{0.9,0.9,0.9}{
            \makecell[{{p{\linewidth}}}]{
                \texttt{\tiny{[GM$|$GM]}}
                \texttt{Maximum number of penalties exceeded} \\
            }
        }
    }
    & & \\ \\

    \theutterance \stepcounter{utterance}  
    & & & \multicolumn{2}{p{0.3\linewidth}}{
        \cellcolor[rgb]{0.9,0.9,0.9}{
            \makecell[{{p{\linewidth}}}]{
                \texttt{\tiny{[GM$|$GM]}}
                \texttt{Initial states:} \\
\texttt{```} \\
\texttt{Player 1:} \\
\texttt{ 1234567} \\
\texttt{╔═╤═════╗} \\
\texttt{║◌│◌A◌◌◌║ 1} \\
\texttt{╟─┤◌T◌◌◌║ 2} \\
\texttt{║O└─┐◌◌◌║ 3} \\
\texttt{║◌◌◌├─┐◌║ 4} \\
\texttt{║◌◌S│◌│◌║ 5} \\
\texttt{╟─┬─┴─┼─╢ 6} \\
\texttt{║P│◌DM│◌║ 7} \\
\texttt{╚═╧═══╧═╝} \\
\\ 
\\ 
\texttt{Player 2:} \\
\texttt{ 1234567} \\
\texttt{╔═╤═════╗} \\
\texttt{║P│S◌D◌◌║ 1} \\
\texttt{╟─┤◌◌◌◌◌║ 2} \\
\texttt{║◌└─┐◌A◌║ 3} \\
\texttt{║◌◌M├─┐◌║ 4} \\
\texttt{║◌T◌│◌│O║ 5} \\
\texttt{╟─┬─┴─┼─╢ 6} \\
\texttt{║◌│◌◌◌│◌║ 7} \\
\texttt{╚═╧═══╧═╝} \\
\\ 
\texttt{```} \\
            }
        }
    }
    & & \\ \\

    \theutterance \stepcounter{utterance}  
    & & & \multicolumn{2}{p{0.3\linewidth}}{
        \cellcolor[rgb]{0.9,0.9,0.9}{
            \makecell[{{p{\linewidth}}}]{
                \texttt{\tiny{[GM$|$GM]}}
                \texttt{End states:} \\
\texttt{```} \\
\texttt{Player 1:} \\
\texttt{ 1234567} \\
\texttt{╔═╤═════╗} \\
\texttt{║◌│◌A◌◌◌║ 1} \\
\texttt{╟─┤◌T◌◌◌║ 2} \\
\texttt{║O└─┐◌◌◌║ 3} \\
\texttt{║◌◌◌├─┐◌║ 4} \\
\texttt{║◌◌S│◌│◌║ 5} \\
\texttt{╟─┬─┴─┼─╢ 6} \\
\texttt{║P│◌DM│◌║ 7} \\
\texttt{╚═╧═══╧═╝} \\
\\ 
\\ 
\texttt{Player 2:} \\
\texttt{ 1234567} \\
\texttt{╔═╤═════╗} \\
\texttt{║P│S◌D◌◌║ 1} \\
\texttt{╟─┤◌◌◌◌◌║ 2} \\
\texttt{║◌└─┐◌A◌║ 3} \\
\texttt{║◌◌M├─┐◌║ 4} \\
\texttt{║◌T◌│◌│O║ 5} \\
\texttt{╟─┬─┴─┼─╢ 6} \\
\texttt{║◌│◌◌◌│◌║ 7} \\
\texttt{╚═╧═══╧═╝} \\
\\ 
\texttt{```} \\
            }
        }
    }
    & & \\ \\

    \theutterance \stepcounter{utterance}  
    & & & \multicolumn{2}{p{0.3\linewidth}}{
        \cellcolor[rgb]{0.9,0.9,0.9}{
            \makecell[{{p{\linewidth}}}]{
                \texttt{\tiny{[GM$|$GM]}}
                \texttt{* success: False} \\
\texttt{* lose: False} \\
\texttt{* aborted: True} \\
\texttt{{-}{-}{-}{-}{-}{-}{-}} \\
\texttt{* Shifts: 0.00} \\
\texttt{* Max Shifts: 12.00} \\
\texttt{* Min Shifts: 6.00} \\
\texttt{* End Distance Sum: 32.45} \\
\texttt{* Init Distance Sum: 32.45} \\
\texttt{* Expected Distance Sum: 29.33} \\
\texttt{* Penalties: 17.00} \\
\texttt{* Max Penalties: 16.00} \\
\texttt{* Rounds: 5.00} \\
\texttt{* Max Rounds: 28.00} \\
\texttt{* Object Count: 7.00} \\
            }
        }
    }
    & & \\ \\

\end{supertabular}
}

\end{document}
